\subsection{交换代数的可分次数}

设 A 是 K 上有限次数 n 的交换代数。对 K 的任意扩张 L，A 到 L 的代数同态的个数 \(h(L)\) 是有限的，并且以 n 为上界（V，第29页，推论）。

引理1. --- 设 \(\Omega\) 是 K 的一个代数闭包；于是对 K 的每个扩张 L，我们有 \(h(L) \leqslant h(\Omega)\) ，当 L 代数闭时取等号。

设 \(L'\) 是 K 在 L 中的代数闭包。对 A 到 L 的每个同态 u，我们有 \([u(A):K] \leqslant n\) ，因此由 V，第18页，命题2，\(u(A) \subset L'\) ；于是我们有 \(h(L') = h(L)\) 。由于 K 的扩张 \(L'\) 同构于 \(\Omega(V, p. 20, Th. 1)\) 的一个子扩张，我们有 \(h(L') \leqslant h(\Omega)\) 。若 L 是代数闭的，则 \(L'\) 是 K 的一个代数闭包；K 的扩张 \(L'\) 与 \(\Omega\) 于是同构（V，第23页，定理2），因此 \(h(L') = h(\Omega)\) ；引理直接由此得出。

根据引理1，数 \(h(\mathbf{L})\) 对所有代数闭扩张 \(\mathbf{L}\) （\(\mathbf{K}\) 的）都有相同的值；这个数将记作 \([\mathbf{A}:\mathbf{K}]_s\) ，并称为 \(\mathbf{A}\) 的可分次数。

设 \(\mathbf{A}\) 与 \(\mathbf{B}\) 是两个有限次数的交换代数（在 \(\mathbf{K}\) 上）。我们将建立公式

\[
[ \mathbf {A} \otimes_ {\mathbf {K}} \mathbf {B}: \mathbf {K} ] _ {s} = [ \mathbf {A}: \mathbf {K} ] _ {s}. [ \mathbf {B}: \mathbf {K} ] _ {s}.\tag{7}
\]

设 \(L\) 是 \(K\) 的一个代数闭扩张，并用 \(\mathcal{H}(A)\) 表示 \(A\) 到 \(L\) 的代数同态之集，类似地定义 \(\mathcal{H}(B)\) 与 \(\mathcal{H}(A \otimes_K B)\) 。根据定义，我们有 Card \(\mathcal{H}(A) = [A : K]_s\) ，以及对 \([B:K]_{s}\) 与 \([A\otimes_{K}B:K]_{s}\) 的相应公式。此外（III，第465页，公式(6)），公式 \((u*v)(a\otimes b)=u(a)v(b)\) 定义了一个双射 \((u,v)\mapsto u*v\) （从 \(\mathcal{H}(A)\times\mathcal{H}(B)\) 到 \(\mathcal{H}(A\otimes_{K}B)\) 上），由此公式(7)成立。

设 \(\mathbf{K}'\) 是 \(\mathbf{K}\) 的一个扩张；我们将证明公式

\[
[ \mathbf {A} _ {(\mathbf {K} ^ {\prime})}: \mathbf {K} ^ {\prime} ] _ {s} = [ \mathbf {A}: \mathbf {K} ] _ {s}.\tag{8}
\]

例如取 L 为 \(K'\) 的一个代数闭包。公式 \(\tilde{u}(x)=u(1\otimes x)\) （对 \(x\in A\)）定义了一个双射 \(u\mapsto\tilde{u}\) ，它在 \(K'\) -同态（由 \(\mathrm{A}_{(\mathrm{K}')}\) 到 L 的）组成的集合与由 A 到 L 的 K-同态组成的集合之间，由此得到 (8)。

最后，假设 \(K'\) 是 K 的有限次扩张；若 \(A'\) 是有限次数的交换 \(K'\) -代数，则它也是有限次数的交换 K-代数，并且我们有 \([A':K] = [A':K']\) . \([K':K]\) （V，第～10页，定理 1）。我们将证明公式

\[
[ \mathrm{A} ^ {\prime}: \mathrm{K} ] _ {s} = [ \mathrm{A} ^ {\prime}: \mathrm{K} ^ {\prime} ] _ {s} \cdot [ \mathrm{K} ^ {\prime}: \mathrm{K} ] _ {s}.\tag{9}
\]

对于设 S（相应地 T）为将 \(K'\) （相应地 \(A'\) ）映入 K 的代数闭包 L 的 K-同态之集；对每个 \(\sigma \in S\) ，我们用 \(T_{\sigma}\) 表示 T 中满足 \(f(\alpha, 1) = \sigma(\alpha)\) （对所有 \(\alpha \in K'\)）的元素 f 的集合。则族 \((\mathrm{T}_{\sigma})_{\sigma \in S}\) 是 T 的一个划分，且我们有 Card \(S = [K': K]_{s}\) ；现在对每个 \(\sigma \in S\) ，集合 \(T_{\sigma}\) 由 \(K'\) -同态组成，它们将 \(A'\) 映入代数闭扩张 \((L, \sigma)\) （\(K'\) 的），因此 Card \(T_{\sigma} = [A': K']_{s}\) ，于是我们已证明了 (9)。

命题4. --- 设A是K上有限次数的交换代数；则 \([A:K]_{s} \leqslant [A:K]\) 当且仅当A是平展时取等号。

设 \(\Omega\) 为 K 的一个代数闭包，且 \(\mathcal{H}\) 为 A 到 \(\Omega\) 的代数同态之集。我们有 Card \(\mathcal{H} = [\mathrm{A}:\mathrm{K}]_s\) ，且 A 是平展的当且仅当 A 被扩张 \(\Omega\) （K 的）对角化（V，第～30页，命题 2）。因此，命题4由 V 第29页的推论得出。

推论1. --- 设 A、B 是 K 上两个有限非零次数的交换代数。则要使代数 \(C = A \otimes_{K} B\) 是平展的，必要且充分的是 A 和 B 都是平展的。

我们有 \([C:K]=[A:K]\) . \([B:K]\) ，以及关于可分次数的相应公式 (7)。进一步，我们有 \([A:K]_{s}\leqslant[A:K]\) 以及关于 B 和 C 的相应公式。由此可得 \([C:K]=[C:K]_{s}\) 当且仅当我们同时有 \([A:K]=[A:K]_{s}\) 与 \([B:K]=[B:K]_{s}\) ；现在只需应用命题 4。

推论2. --- 设K'是K的一个扩域。

\begin{enumerate}
\def\labelenumi{\alph{enumi})}
\item
  要使 K-代数 A 是平展的，必要且充分的是 \(\mathbf{K}'\) -代数 \(\mathbf{A}_{(\mathbf{K}')}\) 是平展的。
\item
  设 \(\mathbf{A}'\) 是 \(\mathbf{K}'\) 上的一个代数，且不约化为 0。要使 \(\mathbf{A}'\) 在 \(\mathbf{K}\) 上是平展的，必要且充分的是 \(\mathbf{A}'\) 在 \(\mathbf{K}'\) 上是平展的，且 \(\mathbf{K}'\) 在 \(\mathbf{K}\) 上是平展的。
\end{enumerate}

我们像推论1那样论证，这次对 \(a\) 应用 (8)，对 \(b\) 应用 (9)。

